\section{Objectifs}
\begin{itemize}
\item Reconnaître fonctions paires et impaires et leur traduction géométrique.
\item Déduire les variations d'une fonction du signe de sa dérivée.
\item Déterminer extrémums, résoudre des optimisations et établir des inégalités.
\item Étudier la position relative de deux courbes.
\end{itemize}

\section{Cours}
\subsection{Parité}
Sur un domaine symétrique par rapport à 0, \texttt{f} est paire si \texttt{f(-x)=f(x)} : la courbe est symétrique par rapport à l'axe des ordonnées. Elle est impaire si \texttt{f(-x)=-f(x)} : la courbe est symétrique par rapport à l'origine.

\subsection{Dérivée et variations}
Si \texttt{f'} est positive sur un intervalle, \texttt{f} y est croissante ; si \texttt{f'} est négative, elle y est décroissante. Si \texttt{f'=0} partout sur un intervalle, \texttt{f} y est constante.

À un extrémum intérieur où \texttt{f} est dérivable, le nombre dérivé est nul. La réciproque est fausse en général : \texttt{f'(a)=0} n'impose pas à lui seul un extrémum.

\subsection{Optimisation et inégalités}
Pour optimiser, on modélise la grandeur par une fonction sur un domaine pertinent, on calcule la dérivée, on établit le tableau de variations puis on interprète l'extrémum.

Pour comparer \texttt{f} et \texttt{g}, on peut étudier le signe de \texttt{f-g}. Les variations peuvent aussi établir une inégalité en localisant un minimum ou un maximum.

\subsection{Second degré et dérivation}
Une fonction \texttt{ax²+bx+c} a une dérivée affine \texttt{2ax+b}. Le signe de cette dérivée permet de retrouver son sens de variation, son extrémum et l'allure de la parabole selon le signe de \texttt{a}.

\section{Algorithmique}
La méthode de Newton peut être explorée dans des cas favorables comme algorithme de recherche d'une racine, sans développer ici la théorie de convergence.

\section{Erreurs fréquentes}
\begin{itemize}
\item Déduire un extrémum de la seule égalité \texttt{f'(a)=0}.
\item Étudier une parité sur un domaine non symétrique.
\item Oublier de restreindre l'optimisation au domaine du problème.
\end{itemize}

\section{Synthèse}
Le signe de la dérivée transforme un calcul local en information globale sur la courbe : variations, extrémums, comparaisons et optimisation.
